\chapter{Hyperparameter Selection and Sensitivity}

The DQN agent's behaviour is governed by a set of hyperparameters whose values were selected based on a combination of established practice in the DRL traffic signal control literature and the specific operational constraints of the Thapathali intersection. Each variable is justified below in terms of its functional role and the rationale for its assigned value.

\section{State Dimensionality (21 dimensions)}

The state vector encodes three categories of intersection information: PCU-weighted load on 7 origin lanes, maximum waiting time on 6 stopline lanes, PCU-weighted load on 5 outgoing lanes, and a 3-dimensional one-hot encoding of the active phase. Each component addresses a distinct informational need. The origin lane PCU values tell the agent where pressure is building; the stopline waiting times capture fairness — whether any single approach has been ignored for too long; the outgoing PCU values reflect whether the intersection is successfully discharging traffic downstream; and the phase encoding prevents the agent from selecting actions without knowing its current operating context. Omitting any of these components would structurally impair the agent's ability to make informed decisions. For instance, an agent without outgoing lane information cannot distinguish between a green phase that is successfully clearing traffic and one that is being blocked by downstream congestion.

\section{Replay Buffer Capacity (10,000 transitions)}

The replay buffer stores past experience tuples to break temporal correlations between consecutive training samples. A buffer that is too small would cause the agent to train predominantly on recent experiences, leading to overfitting to current traffic conditions and instability when demand patterns shift between hours. A buffer that is too large relative to the number of training steps would dilute recent experience with outdated transitions, slowing adaptation. At 10,000 transitions and a 15-second decision step, this buffer spans approximately 41 simulated hours of experience — sufficient to cover the full daily demand cycle multiple times across episodes while remaining computationally tractable. Increasing this to 50,000 would likely improve sample diversity but would require proportionally more training episodes to populate and exploit the buffer effectively.

\section{Batch Size (64 transitions)}

Mini-batch gradient updates use 64 randomly sampled transitions per training step, a value standard in DQN implementations. Smaller batches introduce higher gradient variance, potentially destabilising Q-value updates. Larger batches reduce variance but increase per-step computational cost and can cause the network to overfit to the current batch distribution before the buffer has been adequately repopulated. At 64, each update covers a representative cross-section of stored experiences without excessive computational overhead on CPU-based training.

\section{Discount Factor \( \gamma = 0.95 \)}

The discount factor controls how much the agent values future rewards relative to immediate ones. A value of 1.0 would make the agent treat all future rewards as equally important regardless of how far ahead they occur, which is unsuitable for traffic control where the causal link between a signal decision and its traffic outcome attenuates over time as other vehicles arrive and depart. A value near 0 would make the agent myopic, optimising only the immediate reward of each 15-second step without accounting for downstream queue effects. At \( \gamma = 0.95 \), a reward received five decision steps (75 seconds) ahead is discounted to approximately 77\% of its face value — a horizon long enough to capture the queue-clearing consequences of phase decisions while keeping the training objective stable. Reducing \( \gamma \) to 0.90 would likely cause the agent to be more reactive and less capable of holding a phase strategically to clear a building queue.

\section{Epsilon Start (\( \varepsilon = 1.0 \)), Minimum (\( \varepsilon_{\min} = 0.01 \)), and Decay (\( \varepsilon_{\text{decay}} = 0.985 \))}

The epsilon-greedy strategy governs exploration versus exploitation. Starting at \( \varepsilon = 1.0 \) ensures the agent explores the state space broadly in early episodes rather than immediately exploiting the random initial Q-values of the untrained network, which would bias learning toward the first arbitrary policy encountered. The decay factor of 0.985 applied after each episode causes \( \varepsilon \) to reach approximately 0.18 by episode 100 and 0.01 by episode 300, providing a gradual transition from exploration to exploitation that aligns with the agent's growing ability to make informed decisions as training progresses. A faster decay (e.g., 0.97) would cause the agent to commit to exploitation prematurely, before Q-values have stabilised, risking convergence to a suboptimal policy. A slower decay (e.g., 0.995) would extend random exploration unnecessarily into late training, wasting episodes where the agent has already learned a reasonable policy. The minimum \( \varepsilon = 0.01 \) retains a 1\% probability of random action throughout evaluation to prevent the agent from becoming trapped in a deterministic but suboptimal cycle.

\section{Target Network Update Frequency (every 10 episodes)}

The target network provides stable Q-value targets during training. If the target network updated at every step alongside the Q-network, the training targets would shift continuously, creating a moving-target problem that frequently causes Q-value divergence. Updating too infrequently (e.g., every 50 episodes) would cause the agent to train against outdated targets for too long, particularly during the rapid early learning phase. At 10 episodes, the target network is refreshed frequently enough to reflect meaningful policy improvement while remaining stable enough to provide reliable training signals between updates.

\section{Reward Weighting Factor \( \lambda = 0.0004 \)}

The reward function combines a Max-Pressure throughput term with a PCU-weighted accumulated waiting time penalty scaled by \( \lambda \). The role of \( \lambda \) is to determine how much the agent penalises individual vehicle waiting time relative to the aggregate pressure objective. Without the waiting time term (\( \lambda = 0 \)), the agent could learn a policy that maximises overall throughput while consistently starving one low-volume approach of green time — technically efficient but operationally unacceptable. Too large a \( \lambda \) (e.g., 0.01) would cause the agent to focus disproportionately on eliminating individual vehicle delay, potentially sacrificing overall throughput. The value \( \lambda = 0.0004 \) was selected to keep the wait penalty secondary in magnitude to the pressure term across typical PCU load ranges observed at Thapathali, so that fairness acts as a corrective constraint rather than the primary objective. At typical peak-hour stopline loads of 20–30 PCU and accumulated wait times of 30–60 seconds per vehicle, the wait penalty contributes approximately 0.24–0.72 units per vehicle to the reward, while the pressure term contributes 5–15 units, maintaining the intended hierarchy. Increasing \( \lambda \) by an order of magnitude would be expected to reduce the worst-case wait times on low-volume approaches but at the cost of overall throughput efficiency.

\section{Decision Step Duration (15 seconds)}

Each action is held for 15 seconds of simulation time before the agent re-evaluates the intersection state. This value represents a balance between responsiveness and stability. A shorter step (e.g., 5 seconds) would allow more frequent decisions but would increase the risk of unnecessary phase switching, as queue dynamics at Thapathali do not change meaningfully within a few seconds. Frequent switching also incurs additional yellow-phase transition time (3 seconds each), which would consume a disproportionate share of available green time. A longer step (e.g., 30 seconds) would reduce switching overhead but would limit the agent's ability to respond to rapid changes in approach demand. At 15 seconds, the agent can make approximately 4 decisions per minute — sufficient granularity to track peak-hour demand dynamics without excessive transition overhead. The maximum phase duration of 60 seconds (4 consecutive STAY actions) similarly prevents any single approach from being served indefinitely while still allowing the agent to hold a productive phase through a full queue-clearing cycle.

\section{Yellow Phase Duration (3 seconds)}

A 3-second yellow interval is inserted before every phase transition to reflect safe and realistic signal operation. This value is consistent with standard traffic engineering practice for intersection approach speeds in the 25–40 km/h range typical of urban Kathmandu roads. Reducing this to 1–2 seconds would be unsafe; increasing it to 5 seconds would unnecessarily consume green time, particularly at the 15-second decision step resolution where a 5-second yellow would represent one-third of available action time.

\section{Hidden Layer Architecture (128 \( \rightarrow \) 64 \( \rightarrow \) 32 neurons)}

The three-layer architecture with progressively narrowing hidden layers is designed to extract hierarchical representations from the 21-dimensional state input. The first layer (128 neurons) provides sufficient capacity to capture the full range of pairwise interactions between state features — lane loads, waiting times, and phase indicators — that jointly determine the optimal action. The subsequent layers (64 and 32 neurons) compress this representation into increasingly abstract features, forcing the network to generalise rather than memorise specific state patterns. A deeper or wider network could in principle capture more complex relationships but would require substantially more training data and episodes to converge reliably, which was not feasible within the 300-episode training budget. A shallower network (e.g., single hidden layer) risks underfitting the non-linear structure of traffic state dynamics.

\section{Learning Rate (\( \alpha = 0.001 \)) for Adam Optimiser}

The Adam optimiser's learning rate controls the step size of each gradient update to the Q-network weights. At 0.001, a standard default for Adam in DRL applications, the network updates are large enough to make meaningful progress per batch but small enough to avoid overshooting well-converged Q-value estimates. A higher learning rate (e.g., 0.01) risks Q-value oscillations and divergence, particularly during the early high-variance exploration phase. A lower rate (e.g., 0.0001) would require significantly more training episodes to converge.

\section{Number of Training Episodes (300)}

The training run of 300 episodes was selected based on the epsilon decay schedule: at \( \varepsilon_{\text{decay}} = 0.985 \) per episode, \( \varepsilon \) reaches its minimum of 0.01 by approximately episode 300, meaning the agent transitions fully to exploitation at the same point training concludes. Extending training beyond 300 episodes would primarily refine the policy under full exploitation, offering diminishing returns unless accompanied by renewed exploration. The episode reward trajectory in Figure 6-2 confirms that the moving average stabilises by approximately episode 250, supporting the adequacy of 300 episodes for this intersection configuration.
